\documentclass[11pt,letterpaper]{article}

\usepackage[utf8]{inputenc}
\usepackage[spanish]{babel}
\usepackage[margin=2.5cm]{geometry}
\usepackage{enumitem}
\usepackage{titlesec}
\usepackage{xcolor}
\usepackage{listings}

\definecolor{unicaucaverde}{RGB}{0,90,60}
\titleformat{\section}{\normalfont\Large\bfseries\color{unicaucaverde}}{\thesection}{1em}{}

\lstdefinestyle{cstyle}{
  language=C,
  basicstyle=\ttfamily\small,
  keywordstyle=\color{unicaucaverde}\bfseries,
  commentstyle=\color{gray}\itshape,
  stringstyle=\color{orange!70!black},
  numbers=left,
  numberstyle=\tiny\color{gray},
  frame=single,
  breaklines=true,
  showstringspaces=false,
  tabsize=2
}
\lstset{style=cstyle}

\title{\textcolor{unicaucaverde}{Ejercicios de refuerzo --- Semana 2}\\[4pt]
  \large Condicionales y ciclos en C}
\author{Programación --- Universidad del Cauca}
\date{2026}

\begin{document}
\maketitle

Práctica adicional para consolidar RA2.1 y RA2.2, para trabajo autónomo antes de la Semana 3. Cada
ejercicio incluye el enunciado y una solución completa en C.

\section*{Ejercicio 1 --- Serie de Fibonacci}
Escribe un programa que lea un entero $n$ e imprima los primeros $n$ términos de la serie de
Fibonacci ($0, 1, 1, 2, 3, 5, 8, \dots$, cada término es la suma de los dos anteriores).

\begin{lstlisting}
#include <stdio.h>

int main(void) {
    int n, actual = 0, siguiente = 1, temp;

    printf("Cuantos terminos: ");
    scanf("%d", &n);

    for (int i = 0; i < n; i++) {
        printf("%d ", actual);
        temp = actual + siguiente;
        actual = siguiente;
        siguiente = temp;
    }
    printf("\n");
    return 0;
}
\end{lstlisting}

\section*{Ejercicio 2 --- Suma de dígitos}
Escribe un programa que lea un entero positivo y calcule la suma de sus dígitos, usando \texttt{\%}
y \texttt{/} dentro de un ciclo \texttt{while}.

\begin{lstlisting}
#include <stdio.h>

int main(void) {
    int n, suma = 0;

    printf("Ingrese un numero entero positivo: ");
    scanf("%d", &n);

    while (n > 0) {
        suma += n % 10;   // ultimo digito
        n /= 10;          // se elimina el ultimo digito
    }
    printf("Suma de digitos = %d\n", suma);
    return 0;
}
\end{lstlisting}

\section*{Ejercicio 3 --- Invertir un número}
Escribe un programa que lea un entero positivo y muestre sus dígitos en orden inverso (por
ejemplo, \texttt{1234} $\to$ \texttt{4321}).

\begin{lstlisting}
#include <stdio.h>

int main(void) {
    int n, invertido = 0, digito;

    printf("Ingrese un numero entero positivo: ");
    scanf("%d", &n);

    while (n > 0) {
        digito = n % 10;
        invertido = invertido * 10 + digito;
        n /= 10;
    }
    printf("Numero invertido = %d\n", invertido);
    return 0;
}
\end{lstlisting}

\section*{Ejercicio 4 --- Menú con validación}
Escribe un programa con un menú (1. Sumar, 2. Restar, 3. Salir) usando \texttt{do-while} y
\texttt{switch}, que repita el menú hasta que el usuario elija la opción 3.

\begin{lstlisting}
#include <stdio.h>

int main(void) {
    int opcion;
    float a, b;

    do {
        printf("\n1. Sumar\n2. Restar\n3. Salir\n");
        printf("Elija una opcion: ");
        scanf("%d", &opcion);

        switch (opcion) {
            case 1:
                printf("Ingrese dos numeros: ");
                scanf("%f %f", &a, &b);
                printf("Suma = %.2f\n", a + b);
                break;
            case 2:
                printf("Ingrese dos numeros: ");
                scanf("%f %f", &a, &b);
                printf("Resta = %.2f\n", a - b);
                break;
            case 3:
                printf("Saliendo...\n");
                break;
            default:
                printf("Opcion invalida\n");
        }
    } while (opcion != 3);

    return 0;
}
\end{lstlisting}

\section*{Ejercicio 5 --- Patrón de asteriscos}
Escribe un programa que, usando ciclos anidados, imprima un triángulo de asteriscos de $n$ filas
(fila 1 con 1 asterisco, fila 2 con 2, ..., fila $n$ con $n$).

\begin{lstlisting}
#include <stdio.h>

int main(void) {
    int n;
    printf("Numero de filas: ");
    scanf("%d", &n);

    for (int fila = 1; fila <= n; fila++) {
        for (int col = 1; col <= fila; col++) {
            printf("*");
        }
        printf("\n");
    }
    return 0;
}
\end{lstlisting}
Con $n=4$ imprime:
\begin{verbatim}
*
**
***
****
\end{verbatim}

\section*{Ejercicio 6 --- Clasificación de un dispositivo IoT}
Un dispositivo de instrumentación reporta si tiene sensor, si tiene actuador, su nivel de batería
(0.0 a 100.0) y si está conectado a la red. Escribe un programa que lea estos cuatro datos y
clasifique el dispositivo combinando condiciones con \texttt{\&\&}, \texttt{||} y \texttt{!},
según las reglas siguientes (evaluadas en este orden):
\begin{itemize}[leftmargin=1.2cm]
  \item Sensor \textbf{y} actuador \textbf{y} batería $>50$ \textbf{y} conectado
        $\to$ \emph{IoT avanzado}.
  \item Sensor \textbf{y} batería $\geq 20$ \textbf{y} conectado $\to$ \emph{IoT básico}.
  \item (Sensor \textbf{o} actuador) \textbf{y no} conectado $\to$ \emph{sin conexión}.
  \item Batería $<10$ $\to$ \emph{batería crítica}.
  \item En cualquier otro caso $\to$ \emph{no clasificado}.
\end{itemize}

\begin{lstlisting}
#include <stdio.h>

int main(void) {
    int tieneSensor, tieneActuador, conectado;
    float bateria;

    printf("Tiene sensor? (1=si, 0=no): ");
    scanf("%d", &tieneSensor);
    printf("Tiene actuador? (1=si, 0=no): ");
    scanf("%d", &tieneActuador);
    printf("Nivel de bateria (0.0 a 100.0): ");
    scanf("%f", &bateria);
    printf("Conectado a la red? (1=si, 0=no): ");
    scanf("%d", &conectado);

    if (tieneSensor && tieneActuador && bateria > 50 && conectado) {
        printf("Clasificacion: IoT avanzado\n");
    } else if (tieneSensor && bateria >= 20 && conectado) {
        printf("Clasificacion: IoT basico\n");
    } else if ((tieneSensor || tieneActuador) && !conectado) {
        printf("Clasificacion: sin conexion (revisar red)\n");
    } else if (bateria < 10) {
        printf("Clasificacion: bateria critica, cargar urgente\n");
    } else {
        printf("Clasificacion: no clasificado\n");
    }
    return 0;
}
\end{lstlisting}

\end{document}
